\section{An exact formula for the arch moment}\label{sec:arch}

We now reconnect the polynomial vacuum with a positive path observable.
Cutting a rightward arch at every column between its endpoints counts
it exactly as many times as its horizontal displacement. We use this
pairing of source states to identify the arch moment with a polynomial
coefficient. A stepped-boundary flux calculation then compares its
successive increments with the bridge mass.


For a repeated rhombic column $X=(u_1,\ldots,u_N)$ near a physical
parameter list, let $K_X(k)$ be the bottom-to-bottom partition sum
with displacement $k$, with exactly its two endpoints occupied on
the exterior boundary and one connected path. Set
$m(X)=\sum_{k\ge1}kK_X(k)$. A cut can meet the same arch several
times: either half then contains its source fragment together with
paired fragments, as in Figure~\ref{fig:cut-states}. Using the
south-source column $\mathsf b_N$ from Section~\ref{sec:flux}, put
\begin{equation}\label{eq:01-integrated-source}
 H_N=(I-T_N^{(1)})^{-1}\mathsf b_NP_N,
\end{equation}
and define $H_N^-$ in the reflected strip. Here $P_N$ supplies the
paired fragments before the source column, while the resolvent sums
the columns between that source and the cut. Let $\operatorname{Glue}$
be the bilinear pairing of one-source states that is one when their
occupancies agree and gluing joins the two sources into a single
component with no loops, and zero otherwise. Then
\begin{equation}\label{eq:01-moment-gluing}
 m(X)=\operatorname{Glue}(H_N^-,H_N).
\end{equation}
Indeed expanding the two resolvents places one source on either
side of the cut. For endpoint separation $k$ there are exactly $k$
possible cuts between them. Lemma~\ref{lem:transfer-decay} justifies absolute
convergence of the source sums and the vacuum sums in a neighborhood
of every mixed parameter list. Thus $m$ is rational in the spectral
variables and regular at all such lists.

Related finite-strip current calculations use bilinear link-state
pairings, fusion reductions and interpolation
\cite{DeGierNienhuisPonsaing2010,FeherNienhuis2018}. Those $O(1)$
current formulas use an assumed cancellation of a normalizing factor or additional
symmetry. For the present zero-loop-value arch moment, the source-sector
bound proved below supplies the polynomial control needed for
interpolation.

\begin{proposition}[Arch formula]\label{prop:arch-formula}
Let $\eta=2d-1$ and
\begin{equation}\label{eq:01-g-definition}
 g_X(z)=\frac{p_{N+1}(a_1,\ldots,a_N,z)}{p_N(a_1,\ldots,a_N)}.
\end{equation}
Then $g_X$ is monic of degree $N$ for generic $X$, and
\begin{equation}\label{eq:01-arch-formula}
 m(X)=\sum_{i=1}^Na_i-\eta[z^{N-1}]g_X(z).
\end{equation}
The right side is understood by continuation at removable
singularities.
\end{proposition}

\begin{proof}
We identify $m$ from its row reductions and a separate degree bound
of $2N$ for $p_N^2m$ in each $a_i$. The source-sector limit supplies
that bound; the scalar reductions then determine the coefficient in
\eqref{eq:01-arch-formula}.

\paragraph{Symmetries and reductions.}
Moving an exchange box through both halves of
\eqref{eq:01-moment-gluing}, using \eqref{eq:transfer-intertwiners},
proves symmetry in the rows. At an ascending merge or deletion
specialization, the same identities for the vacuum and south-source
columns give
\[
 P_N=S_sP_{\mathrm{red}},\qquad H_N=S_sH_{\mathrm{red}}.
\]
The reflected column identity \eqref{eq:column-reflection} gives
\[
 S_s^{\mathsf t}P_N^-=P_{\mathrm{red}}^-,\qquad
 S_s^{\mathsf t}H_N^-=H_{\mathrm{red}}^-.
\]
For the first reflected equality, the projected vector is stationary
with empty coordinate one. For the second, the projected source
equation is
\begin{equation}\label{eq:arch-reflected-source}
 (I-T_{\rm red}^{(1),-})S_s^{\mathsf t}H_N^-
   =\mathsf b_{\rm red}^-S_s^{\mathsf t}P_N^-
   =\mathsf b_{\rm red}^-P_{\rm red}^-.
\end{equation}
Uniqueness near a mixed specialization and rational continuation
therefore identify the projected vector with $H_{\rm red}^-$.
Transposing the splitter inside the gluing pairing now reduces $m$
to the smaller strip. Thus the two halves reduce by insertion and
projection, respectively; no identity $S_s^{\mathsf t}S_s=I$ is
required. Occupancy signs cancel under a $\pi$ shift of a parameter.

There is also a separate reflection symmetry
$u_j\leftrightarrow3\lambda-u_j$. By exchange, put that row on top.
Since all top ports are vacant, each visit to the row is an interval
that turns up and then back down. An interval with displacement
$\ell\ge1$ has weight $AUB^{\ell-1}$, which is invariant under this
reflection. The intervals are ordered and disjoint. The argument is
valid for absolutely convergent sums near physical parameters and
hence as a rational identity. If $a_j=c$, then $AU=0$ and these
excursions vanish, so the flat row can simply be omitted.

\paragraph{A polynomial bound from the source sector.}
To obtain the degree bound, first use a north source. Add a last row
$y$ to the vacuum strip, and identify its occupied last-slot sector
with one-source states on the earlier rows. Last-box stationarity is
\begin{equation}\label{eq:01-north-source-limit}
 (P_{N+1})_o
 =B(y)T_N^{(1)}(P_{N+1})_o
   +U(y)\mathsf b_N^{\rm north}(P_{N+1})_v.
\end{equation}
As $b=a(y)\to\infty$, Lemmas~\ref{lem:scalar-reductions} and
\ref{lem:vacuum-degree} give $(P_{N+1})_v\to P_N$, while $B(y)\to1$.
Thus $(P_{N+1})_o/U(y)$ tends to the integrated north-source vector.
Since $U(y)=2d\sin y/(b+C)$, its numerator after multiplication by
$p_N$ is exactly
\[
 \frac1{2d}[b^{N-1}]\frac{(W_{N+1})_o}{\sin y}.
\]
The lemma bounds its Laurent exponents in each earlier variable by
$2N$. Reflection from top to bottom reverses row order and replaces
$u_i$ by $3\lambda-u_i$; it preserves $p_N$ and these exponent bounds.
Hence both $p_NH_N$ and $p_NH_N^-$ satisfy the same bound.
By \eqref{eq:01-moment-gluing}, $p_N^2m$ has exponent bound $4N$.
Its separate reflection and $\pi$-periodicity imply that it is a
polynomial in $a_i=\cos(2u_i-3\lambda)$ of degree at most $2N$ in
each variable. Indeed, a $\pi$-periodic Laurent polynomial is a
Laurent polynomial in $z=e^{2iu_i}$, and invariance under
$z\mapsto e^{6i\lambda}/z$ makes it a polynomial in the corresponding
scaled sum $a_i$. Row exchange makes this polynomial symmetric.

\paragraph{Interpolation.}
Let the right side of \eqref{eq:01-arch-formula} be $\widetilde m$.
The scalar identities give the exact quotient rules
\[
 g_{\rm merge}(z)=(z+f_0)g_{\rm red}(z),\quad
 g_{\rm delete}(z)=(z-f_{-6})(z-f_{12})g_{\rm red}(z),\quad
 g_{\rm flat}(z)=(z+C)g_{\rm red}(z).
\]
The elementary identities
\[
 f_{-2}+f_2-f_0=\eta f_0,\qquad
 f_0+f_6=-\eta(f_{-6}+f_{12}),\qquad c=\eta C
\]
show that $\widetilde m$ has the same three reductions as $m$.
Both $p_N^2m$ and $p_N^2\widetilde m$ have degree at most $2N$ in each
$a_i$. Induct on $N$. In one distinguished variable, every other
argument $a_j=\cos\theta_j$ supplies the two merge values
$\cos(\theta_j\pm4\lambda)$ and the two deletion values
$\cos(\theta_j\pm6\lambda)$, in addition to the flat value $c$.
These are $4N-3$ distinct roots for generic other arguments, more
than $2N$ for every $N\ge2$. Therefore the polynomial difference
vanishes. The initial cases are $m(\varnothing)=0$ and
\[
 m(u)=\frac{A(u)U(u)}{(1-B(u))^2}
     =\frac{a(u)-c}{1+d}
     =a(u)-\eta(a(u)+\rho).
\]
In the size-zero quotient convention use $g_0=1$ and
$[z^{-1}]g_0=0$. This completes the induction.
\end{proof}

At $X=(\lambda,\ldots,\lambda)$, each rhombic row splits into the
two triangular tiles of one band, with weight $\rho$ per visited
vertex. Thus $a_i=C$ and $m(X)=m_N$ in the port convention of the
introduction. In particular,
\begin{equation}\label{eq:physical-arch-coefficient}
 m_N=NC-\eta\lim_{a_1,\ldots,a_N\to C}
                [z^{N-1}]g_X(z).
\end{equation}
This coefficient has a removable limit because $m$ is regular at the
physical list. Section~\ref{sec:integral} will prove that the entire
monic polynomial $g_X$ has a unique limit and will evaluate this
coefficient through a positive integral.

\begin{lemma}[Moment increments and bridge mass]\label{lem:moment-increments}
For homogeneous honeycomb strips,
$c\mathcal A_N+\mathcal B_N=1$, the bridge mass is nonincreasing,
and $m_{N+1}-m_N\asymp\mathcal B_N$ uniformly in $N\ge1$.
\end{lemma}

\begin{proof}
Apply Lemma~\ref{lem:boundary-flux} to parallelogram truncations of the
strip. The left and right boundary terms tend to zero by
Lemma~\ref{lem:transfer-decay}: after the source column, propagate in the
one-source sector to the distant cut, where the exit is the only
occupied boundary port. The remaining vacuum states are bounded.
The other sums increase to their infinite-strip limits. Bottom exits
have winding $\pm\pi$ and top exits winding zero, so the real part is
$c\mathcal A_N+\mathcal B_N=1$. Increasing strip height increases the
arch sum by inclusion, hence decreases $\mathcal B_N$.

Add one triangular row above the right half of $\mathcal S_N$, in
columns $1,2,\ldots$, creating a stepped strip with a single side
port $p$ at the left end of the new row. Denote its bottom arch
kernel by $\overline K(i,j)$; see Figure~\ref{fig:stepped-strip}.

\begin{figure}[htbp]
\centering
\begin{tikzpicture}[x=1cm,y=1cm,font=\small,>=Latex]
 \fill[blue!4] (-3.2,0) rectangle (3.55,1.7);
 \fill[orange!12] (0,1.7) rectangle (3.55,2.5);
 \draw[thick] (-3.2,0)--(3.55,0);
 \draw[thick] (-3.2,1.7)--(0,1.7)--(0,2.5)--(3.55,2.5);
 \draw[dashed] (0,1.7)--(3.55,1.7);
 \draw[black!50,dashed] (-2.6,-.15)--(-2.6,1.95);
 \draw[black!50,dashed] (3.45,-.15)--(3.45,2.7);
 \node[above] at (-2.6,1.95) {$-L$};
 \node[above] at (3.45,2.7) {$L$};
 \draw[fill=white,thick] (0,2.1) circle (2pt);
 \node[left] at (0,2.1) {$p$};
 \node at (-1.0,.35) {base strip: $N$ bands};
 \node at (1.65,2.15) {one added band};
 \foreach \x/\j in {.5/1,1.1/2,1.7/3,2.3/4,2.9/5}{
  \fill (\x,1.7) circle (1.5pt);
  \node at (\x,.9) {$\j$};
 }
 \draw[very thick,blue!65!black] (0,2.1) .. controls (.5,2.1) and (.5,2.0) .. (.5,1.7);
 \draw[very thick,blue!65!black] (.5,1.7) .. controls (.5,1.0) and (1.1,1.0) .. (1.1,1.7);
 \draw[very thick,blue!65!black] (1.1,1.7) .. controls (1.1,2.0) and (1.7,2.0) .. (1.7,1.7);
 \draw[very thick,blue!65!black] (1.7,1.7) .. controls (1.7,1.0) and (2.3,1.0) .. (2.3,1.7);
 \draw[very thick,blue!65!black] (2.3,1.7) .. controls (2.3,2.0) and (2.9,2.0) .. (2.9,1.7);
 \draw[very thick,blue!65!black] (2.9,1.7) .. controls (3.3,1.3) and (3.3,.7) .. (3.3,0);
 \draw[fill=white,thick] (3.3,0) circle (2pt);
\end{tikzpicture}
\caption{The stepped strip used to compare $m_{N+1}$ and $m_N$.
The drawing is schematic: curves represent path pieces, not individual
lattice edges. Interface visits are numbered from left to right.
A single connected path forces alternating consecutive pairings in the
added row and the base strip; the final visit connects to the bottom.
Moving the two dashed cuts to infinity recovers the two homogeneous
arch moments.}
\label{fig:stepped-strip}
\end{figure}


From a bottom source $i$, the imaginary
part of the boundary identity is
\begin{equation}\label{eq:01-step-identity}
 \sin(3\lambda)\left(
   \sum_{j<i}\overline K(i,j)-\sum_{j>i}\overline K(i,j)
 \right)+\sin\lambda\,Z(i,p)=0.
\end{equation}
Indeed horizontal top exits have winding zero, and closing along the
leftward boundary gives winding $\pi/3$ at $p$. Lateral truncation
errors vanish as before, dominated by a strip of height $N+1$.

Sum \eqref{eq:01-step-identity} over $-L\le i\le L$. Reversal pairs
and cancels all internal arches. The uncanceled terms cross the two
ends of this interval. After translating each end to the origin,
the stepped domains converge locally to strips of heights $N$ and
$N+1$, respectively. All crossing sums are dominated by the first
moment in the strip of height $N+1$, which is finite by transfer
decay. This also controls paths leaving an arbitrarily large lateral
window for each fixed pair of ports. Dominated convergence therefore
gives the exact identity
\begin{equation}\label{eq:01-step-moment}
 m_{N+1}-m_N=\frac{\sin\lambda}{\sin(3\lambda)}
                 \sum_i Z(p,i).
\end{equation}

Turning down immediately from $p$ and following an arbitrary bridge
in the base strip yields
$\sum_iZ(p,i)\ge\rho\mathcal B_N$. The top-source bridge mass in the
base strip equals its bottom-source mass by translation and reversal.

For the upper bound, order the path's visits to the interface between
the base strip and the new row from left to right. In the new row,
whose top is vacant, disjoint intervals pair $p$ with the first visit
and then pair subsequent visits consecutively. The base portions are
a noncrossing pairing with one unpaired visit connected to the bottom;
that defect cannot be enclosed. Since the union is a single path,
the first base visit must pair with its next neighbor: pairing it
farther away would leave a closed component inside. Removing these
two visits and repeating proves that base visits also pair
consecutively, with the last visit the defect. Thus the path encounters
the interface in increasing horizontal order.

At $u=\lambda$, the initial new-row segment has total weight at most
$A/(1-B)$. A one-sided base return arch has total mass at most
$1/(2c)$, using the boundary identity from the top and equality of
the two directions by translation and reversal. A one-sided return
through the new row has total mass
$AU/(1-B)=2cB$. The final base segment has total mass
$\mathcal B_N$. Ignoring avoidance between pieces only increases the
sum, since the base weights are products over vertices. Therefore
\[
 \sum_iZ(p,i)
 \le\frac{A}{1-B}\sum_{r\ge0}
       \left(\frac1{2c}\,2cB\right)^r\mathcal B_N
 =\frac{A}{(1-B)^2}\mathcal B_N.
\]
Here $B=B(\lambda)=\rho^2<1$, so the constants are independent of
$N$. Combine this with \eqref{eq:01-step-moment}.
\end{proof}
